\documentclass[12pt,a4paper]{article}

% Paketler
\usepackage[utf8]{inputenc}
\usepackage[turkish]{babel}
\usepackage{graphicx}
\usepackage{xcolor}
\usepackage{listings}
\usepackage{geometry}
\usepackage{fancyhdr}
\usepackage{titlesec}
\usepackage{tcolorbox}
\usepackage{hyperref}
\usepackage{longtable}
\usepackage{booktabs}
\usepackage{array}
\usepackage{enumitem}
\usepackage{tikz}
\usetikzlibrary{shapes.geometric, arrows, positioning, shadows}

% Sayfa ayarları
\geometry{
    a4paper,
    left=2.5cm,
    right=2.5cm,
    top=3cm,
    bottom=3cm
}

% Başlık stili
\titleformat{\section}
{\color{blue!70!black}\normalfont\Large\bfseries}
{\thesection}{1em}{}

\titleformat{\subsection}
{\color{blue!50!black}\normalfont\large\bfseries}
{\thesubsection}{1em}{}

% Header ve Footer
\pagestyle{fancy}
\fancyhf{}
\fancyhead[L]{\textit{Tickly Mobile}}
\fancyhead[R]{\textit{Flutter iOS Uygulaması}}
\fancyfoot[C]{\thepage}

% Hyperlink ayarları
\hypersetup{
    colorlinks=true,
    linkcolor=blue,
    filecolor=magenta,      
    urlcolor=cyan,
    pdftitle={Tickly Mobile Dokümantasyonu},
    pdfpagemode=FullScreen,
}

% Kod bloğu ayarları (Dart için)
\lstdefinelanguage{Dart}{
    keywords={abstract, as, assert, async, await, break, case, catch, class, const, continue, default, deferred, do, dynamic, else, enum, export, extends, external, factory, false, final, finally, for, get, if, implements, import, in, is, library, new, null, operator, part, rethrow, return, set, static, super, switch, this, throw, true, try, typedef, var, void, while, with, yield},
    keywordstyle=\color{blue}\bfseries,
    ndkeywords={String, int, double, bool, List, Map, Future, Stream, Widget, State, StatelessWidget, StatefulWidget, BuildContext, Provider, TextEditingController},
    ndkeywordstyle=\color{purple}\bfseries,
    sensitive=true,
    comment=[l]{//},
    morecomment=[s]{/*}{*/},
    commentstyle=\color{green!60!black}\ttfamily,
    stringstyle=\color{red}\ttfamily,
    morestring=[b]',
    morestring=[b]"
}

\lstset{
    basicstyle=\ttfamily\footnotesize,
    breaklines=true,
    frame=single,
    backgroundcolor=\color{gray!10},
    numbers=left,
    numberstyle=\tiny\color{gray},
    language=Dart
}

% TikZ stil tanımları
\tikzstyle{screen} = [rectangle, rounded corners, minimum width=3cm, minimum height=1.2cm, text centered, draw=black, fill=blue!20, drop shadow]
\tikzstyle{service} = [rectangle, minimum width=3cm, minimum height=1cm, text centered, draw=black, fill=green!20]
\tikzstyle{arrow} = [thick,->,>=stealth]

\begin{document}

% Kapak Sayfası
\begin{titlepage}
    \centering
    \vspace*{2cm}
    
    {\Huge \textbf{TICKLY MOBILE}}\\[0.5cm]
    {\LARGE Flutter iOS Uygulaması}\\[2cm]
    
    {\Large \textbf{Mobil Uygulama Dokümantasyonu}}\\[1cm]
    
    \vfill
    
    \begin{tcolorbox}[colback=blue!5!white,colframe=blue!75!black,title=Proje Bilgileri]
        \begin{itemize}[leftmargin=*]
            \item \textbf{Proje Adı:} Tickly Mobile - iOS App
            \item \textbf{Platform:} iOS (iPhone/iPad)
            \item \textbf{Framework:} Flutter 3.16+
            \item \textbf{Dil:} Dart 3.2+
            \item \textbf{Versiyon:} 1.0.0
            \item \textbf{Minimum iOS:} 12.0
            \item \textbf{Tarih:} Ocak 2026
        \end{itemize}
    \end{tcolorbox}
    
    \vfill
    
    {\large \today}
\end{titlepage}

% İçindekiler
\tableofcontents
\newpage

% Giriş
\section{Proje Genel Bakış}

\subsection{Uygulama Tanımı}

\textbf{Tickly Mobile}, Tickly Help Desk sisteminin iOS mobil uygulamasıdır. Flutter framework'ü kullanılarak geliştirilmiştir ve kullanıcıların mobil cihazlarından ticket oluşturmasına, takip etmesine ve yönetmesine olanak sağlar.

\begin{tcolorbox}[colback=green!5!white,colframe=green!75!black,title=Ana Hedef]
Kullanıcıların her yerden, her zaman ticket sistemine erişim sağlaması, anında bildirim alması ve mobil-friendly bir deneyim yaşaması.
\end{tcolorbox}

\subsection{Ana Özellikler}

\begin{itemize}
    \item \textbf{Kullanıcı Kimlik Doğrulama:} JWT token tabanlı güvenli giriş
    \item \textbf{Ticket Listeleme:} Tüm ticket'ları görüntüleme ve filtreleme
    \item \textbf{Ticket Oluşturma:} Yeni ticket açma, dosya ekleme
    \item \textbf{Ticket Detayı:} Ticket bilgileri, yorumlar, geçmiş
    \item \textbf{Yorum Sistemi:} Ticket'lara yorum ekleme
    \item \textbf{Dosya Ekleme:} Fotoğraf ve doküman yükleme
    \item \textbf{Push Notifications:} Anlık bildirimler (opsiyonel)
    \item \textbf{Dashboard:} İstatistikler ve grafikler
    \item \textbf{Dark Mode:} Koyu tema desteği
    \item \textbf{Offline Support:} Temel veri önbelleği
\end{itemize}

\subsection{Desteklenen Platformlar}

\begin{table}[h]
\centering
\begin{tabular}{|l|c|}
\hline
\textbf{Platform} & \textbf{Minimum Versiyon} \\
\hline
iOS & 12.0+ \\
iPadOS & 12.0+ \\
\hline
\end{tabular}
\caption{Platform Gereksinimleri}
\end{table}

\textbf{Not:} Uygulama öncelikli olarak iOS için geliştirilmiştir, ancak Flutter'ın cross-platform doğası sayesinde Android'e de kolayca port edilebilir.

\newpage

\section{Teknoloji Stack}

\subsection{Core Technologies}

\begin{longtable}{|l|c|p{7cm}|}
\hline
\textbf{Teknoloji} & \textbf{Versiyon} & \textbf{Kullanım Amacı} \\
\hline
\endfirsthead
\hline
\textbf{Teknoloji} & \textbf{Versiyon} & \textbf{Kullanım Amacı} \\
\hline
\endhead
Flutter & 3.16+ & UI framework, cross-platform development \\
\hline
Dart & 3.2+ & Programlama dili, null-safety \\
\hline
Material Design & 3.0 & UI component library \\
\hline
Cupertino Widgets & - & iOS-native görünüm \\
\hline
\caption{Core Teknolojiler}
\end{longtable}

\subsection{Flutter Packages}

\begin{longtable}{|l|c|p{7cm}|}
\hline
\textbf{Package} & \textbf{Versiyon} & \textbf{Kullanım Amacı} \\
\hline
\endfirsthead
\hline
\textbf{Package} & \textbf{Versiyon} & \textbf{Kullanım Amacı} \\
\hline
\endhead
provider & 6.1.1 & State management, dependency injection \\
\hline
http & 1.1.0 & HTTP client, REST API çağrıları \\
\hline
shared\_preferences & 2.2.2 & Local storage, token saklama \\
\hline
image\_picker & 1.0.7 & Kamera ve galeri erişimi \\
\hline
file\_picker & 6.1.1 & Dosya seçme, doküman yükleme \\
\hline
signalr\_netcore & 1.3.7 & Real-time WebSocket bağlantısı \\
\hline
fl\_chart & 0.66.0 & Grafikler ve chart'lar \\
\hline
flutter\_local\_notifications & 16.3.0 & Push notification desteği \\
\hline
intl & 0.18.1 & Internationalization, tarih formatı \\
\hline
\caption{Flutter Packages}
\end{longtable}

\subsection{Backend Integration}

\textbf{API Base URL:}
\begin{itemize}
    \item Development: \texttt{http://localhost:5000/api}
    \item Production: \texttt{https://api.tickly.yourcompany.com/api}
\end{itemize}

\textbf{Kullanılan Endpoint'ler:}
\begin{itemize}
    \item \texttt{POST /api/auth/login} - Kullanıcı girişi
    \item \texttt{GET /api/auth/me} - Kullanıcı bilgisi
    \item \texttt{GET /api/tickets} - Ticket listesi
    \item \texttt{POST /api/tickets} - Yeni ticket oluştur
    \item \texttt{GET /api/tickets/\{id\}} - Ticket detayı
    \item \texttt{POST /api/tickets/\{id\}/comments} - Yorum ekle
    \item \texttt{PUT /api/tickets/\{id\}/status} - Durum güncelle
    \item \texttt{GET /api/departments} - Departman listesi
    \item \texttt{GET /api/categories} - Kategori listesi
\end{itemize}

\newpage

\section{Uygulama Mimarisi}

\subsection{Genel Mimari}

Tickly Mobile, \textbf{Provider} pattern kullanarak state management sağlar ve katmanlı mimari prensiplerine uygun geliştirilmiştir.

\begin{figure}[h]
\centering
\begin{tikzpicture}[node distance=2.5cm]

% UI Layer
\node (ui) [screen, fill=blue!30] {UI LAYER\\Screens \& Widgets};

% Provider Layer
\node (provider) [service, below of=ui, fill=green!30] {PROVIDER LAYER\\State Management};

% Service Layer
\node (service) [service, below of=provider, fill=yellow!30] {SERVICE LAYER\\API \& Business Logic};

% Model Layer
\node (model) [service, below of=service, fill=orange!30] {MODEL LAYER\\Data Models};

% Backend
\node (backend) [service, below of=model, fill=red!30] {BACKEND API\\REST + SignalR};

% Arrows
\draw [arrow, <->] (ui) -- (provider);
\draw [arrow, <->] (provider) -- (service);
\draw [arrow, <->] (service) -- (model);
\draw [arrow, <->] (model) -- node[anchor=east] {HTTP/WebSocket} (backend);

\end{tikzpicture}
\caption{Tickly Mobile Katmanlı Mimari}
\end{figure}

\subsection{Proje Klasör Yapısı}

\begin{lstlisting}[language=bash, numbers=none]
tickly_mobile/
├── lib/
│   ├── main.dart                    # Entry point
│   │
│   ├── models/                      # Data Models
│   │   ├── user.dart
│   │   ├── ticket.dart
│   │   ├── comment.dart
│   │   ├── department.dart
│   │   ├── category.dart
│   │   └── attachment.dart
│   │
│   ├── services/                    # Business Logic & API
│   │   ├── api_service.dart        # Base HTTP client
│   │   ├── auth_service.dart       # Authentication
│   │   ├── ticket_service.dart     # Ticket operations
│   │   ├── notification_service.dart# Push notifications
│   │   └── signalr_service.dart    # Real-time messaging
│   │
│   ├── providers/                   # State Management
│   │   ├── auth_provider.dart      # Auth state
│   │   ├── ticket_provider.dart    # Ticket state
│   │   └── theme_provider.dart     # Theme state
│   │
│   ├── screens/                     # UI Screens
│   │   ├── login_screen.dart
│   │   ├── ticket_list_screen.dart
│   │   ├── create_ticket_screen.dart
│   │   ├── ticket_detail_screen.dart
│   │   ├── dashboard_screen.dart
│   │   └── profile_screen.dart
│   │
│   ├── widgets/                     # Reusable Widgets
│   │   ├── ticket_card.dart
│   │   ├── comment_item.dart
│   │   ├── custom_button.dart
│   │   └── loading_indicator.dart
│   │
│   └── utils/                       # Utilities
│       ├── constants.dart
│       ├── theme.dart
│       └── validators.dart
│
├── pubspec.yaml                     # Dependencies
├── ios/                             # iOS-specific files
└── test/                            # Unit tests
\end{lstlisting}

\newpage

\section{Ekran Tasarımları ve Akışı}

\subsection{Uygulama Ekranları}

\begin{enumerate}
    \item \textbf{Splash Screen}
    \begin{itemize}
        \item Uygulama açılış ekranı
        \item Token kontrolü yapar
        \item Authenticated ise Ticket List'e, değilse Login'e yönlendirir
    \end{itemize}
    
    \item \textbf{Login Screen}
    \begin{itemize}
        \item Email ve şifre ile giriş
        \item "Beni Hatırla" seçeneği
        \item JWT token alır ve shared\_preferences'a kaydeder
    \end{itemize}
    
    \item \textbf{Ticket List Screen}
    \begin{itemize}
        \item Tüm ticket'ları listeler
        \item Filtreleme: Durum, öncelik, departman
        \item Arama fonksiyonu
        \item Pull-to-refresh
        \item FAB (Floating Action Button) ile yeni ticket
    \end{itemize}
    
    \item \textbf{Create Ticket Screen}
    \begin{itemize}
        \item Başlık, açıklama, departman, kategori seçimi
        \item Öncelik seçimi (Low, Medium, High, Critical)
        \item Dosya ve fotoğraf ekleme
        \item Form validation
    \end{itemize}
    
    \item \textbf{Ticket Detail Screen}
    \begin{itemize}
        \item Ticket bilgileri (başlık, durum, öncelik, atanan kişi)
        \item Yorumlar listesi
        \item Yeni yorum ekleme
        \item Durum güncelleme (Agent ve Admin için)
        \item Ekleri görüntüleme
    \end{itemize}
    
    \item \textbf{Dashboard Screen}
    \begin{itemize}
        \item Özet istatistikler (açık, çözülen, kapalı ticket sayıları)
        \item Grafikler (fl\_chart ile)
        \item Son aktiviteler
    \end{itemize}
    
    \item \textbf{Profile Screen}
    \begin{itemize}
        \item Kullanıcı bilgileri
        \item Tema değiştirme (Light/Dark)
        \item Bildirim ayarları
        \item Çıkış yapma
    \end{itemize}
\end{enumerate}

\subsection{Navigasyon Akışı}

\begin{figure}[h]
\centering
\begin{tikzpicture}[node distance=2cm]

\node (splash) [screen] {Splash Screen};
\node (login) [screen, below left of=splash, xshift=-1cm] {Login Screen};
\node (tickets) [screen, below right of=splash, xshift=1cm] {Ticket List};
\node (create) [screen, below of=tickets, yshift=-1cm] {Create Ticket};
\node (detail) [screen, right of=create, xshift=2cm] {Ticket Detail};
\node (dashboard) [screen, left of=create, xshift=-2cm] {Dashboard};

\draw [arrow] (splash) -- node[anchor=south east] {Not Auth} (login);
\draw [arrow] (splash) -- node[anchor=south west] {Auth} (tickets);
\draw [arrow] (login) -- node[anchor=west] {Login Success} (tickets);
\draw [arrow] (tickets) -- node[anchor=east] {FAB} (create);
\draw [arrow] (tickets) -- node[anchor=north] {Tap Ticket} (detail);
\draw [arrow] (tickets) -- node[anchor=north] {Bottom Nav} (dashboard);

\end{tikzpicture}
\caption{Navigasyon Akış Diyagramı}
\end{figure}

\newpage

\section{State Management ve Provider Pattern}

\subsection{Provider Mimarisi}

Tickly Mobile, state management için \textbf{Provider} package'ını kullanır. Her major özellik için ayrı provider'lar oluşturulmuştur.

\subsubsection{AuthProvider}

\begin{lstlisting}[language=Dart]
class AuthProvider extends ChangeNotifier {
  String? _token;
  User? _user;
  bool _isLoading = false;
  
  bool get isAuthenticated => _token != null;
  User? get user => _user;
  bool get isLoading => _isLoading;
  
  Future<bool> login(String email, String password) async {
    _isLoading = true;
    notifyListeners();
    
    try {
      final response = await AuthService.login(email, password);
      _token = response['token'];
      _user = User.fromJson(response['user']);
      
      // Save token to SharedPreferences
      await _saveToken(_token!);
      
      notifyListeners();
      return true;
    } catch (e) {
      _isLoading = false;
      notifyListeners();
      return false;
    }
  }
  
  Future<void> logout() async {
    _token = null;
    _user = null;
    await _clearToken();
    notifyListeners();
  }
  
  Future<void> _saveToken(String token) async {
    final prefs = await SharedPreferences.getInstance();
    await prefs.setString('auth_token', token);
  }
  
  Future<void> _clearToken() async {
    final prefs = await SharedPreferences.getInstance();
    await prefs.remove('auth_token');
  }
}
\end{lstlisting}

\subsubsection{TicketProvider}

\begin{lstlisting}[language=Dart]
class TicketProvider extends ChangeNotifier {
  List<Ticket> _tickets = [];
  Ticket? _currentTicket;
  bool _isLoading = false;
  String? _error;
  
  List<Ticket> get tickets => _tickets;
  Ticket? get currentTicket => _currentTicket;
  bool get isLoading => _isLoading;
  String? get error => _error;
  
  Future<void> fetchTickets() async {
    _isLoading = true;
    _error = null;
    notifyListeners();
    
    try {
      _tickets = await TicketService.getTickets();
      _isLoading = false;
      notifyListeners();
    } catch (e) {
      _error = e.toString();
      _isLoading = false;
      notifyListeners();
    }
  }
  
  Future<bool> createTicket(TicketCreateDto dto) async {
    try {
      final newTicket = await TicketService.createTicket(dto);
      _tickets.insert(0, newTicket);
      notifyListeners();
      return true;
    } catch (e) {
      return false;
    }
  }
  
  Future<void> fetchTicketDetail(int id) async {
    _isLoading = true;
    notifyListeners();
    
    try {
      _currentTicket = await TicketService.getTicketById(id);
      _isLoading = false;
      notifyListeners();
    } catch (e) {
      _error = e.toString();
      _isLoading = false;
      notifyListeners();
    }
  }
}
\end{lstlisting}

\newpage

\section{API Entegrasyonu}

\subsection{Base API Service}

\begin{lstlisting}[language=Dart]
class ApiService {
  static const String baseUrl = 'http://localhost:5000/api';
  
  static Future<Map<String, dynamic>> get(
    String endpoint,
    {String? token}
  ) async {
    final url = Uri.parse('$baseUrl$endpoint');
    final headers = {
      'Content-Type': 'application/json',
      if (token != null) 'Authorization': 'Bearer $token',
    };
    
    final response = await http.get(url, headers: headers);
    
    if (response.statusCode == 200) {
      return json.decode(response.body);
    } else {
      throw Exception('Failed to load data');
    }
  }
  
  static Future<Map<String, dynamic>> post(
    String endpoint,
    Map<String, dynamic> body,
    {String? token}
  ) async {
    final url = Uri.parse('$baseUrl$endpoint');
    final headers = {
      'Content-Type': 'application/json',
      if (token != null) 'Authorization': 'Bearer $token',
    };
    
    final response = await http.post(
      url,
      headers: headers,
      body: json.encode(body),
    );
    
    if (response.statusCode == 200 || 
        response.statusCode == 201) {
      return json.decode(response.body);
    } else {
      throw Exception('Failed to post data');
    }
  }
}
\end{lstlisting}

\subsection{Ticket Service}

\begin{lstlisting}[language=Dart]
class TicketService {
  static Future<List<Ticket>> getTickets() async {
    final token = await _getToken();
    final data = await ApiService.get('/tickets', token: token);
    
    return (data as List)
        .map((json) => Ticket.fromJson(json))
        .toList();
  }
  
  static Future<Ticket> getTicketById(int id) async {
    final token = await _getToken();
    final data = await ApiService.get(
      '/tickets/$id',
      token: token
    );
    
    return Ticket.fromJson(data);
  }
  
  static Future<Ticket> createTicket(
    TicketCreateDto dto
  ) async {
    final token = await _getToken();
    final data = await ApiService.post(
      '/tickets',
      dto.toJson(),
      token: token,
    );
    
    return Ticket.fromJson(data);
  }
  
  static Future<String?> _getToken() async {
    final prefs = await SharedPreferences.getInstance();
    return prefs.getString('auth_token');
  }
}
\end{lstlisting}

\newpage

\section{Data Models}

\subsection{Ticket Model}

\begin{lstlisting}[language=Dart]
class Ticket {
  final int id;
  final String title;
  final String description;
  final TicketStatus status;
  final TicketPriority priority;
  final int departmentId;
  final String departmentName;
  final int? categoryId;
  final String? categoryName;
  final String creatorId;
  final String creatorName;
  final String? assignedToId;
  final String? assignedToName;
  final DateTime createdAt;
  final DateTime? updatedAt;
  
  Ticket({
    required this.id,
    required this.title,
    required this.description,
    required this.status,
    required this.priority,
    required this.departmentId,
    required this.departmentName,
    this.categoryId,
    this.categoryName,
    required this.creatorId,
    required this.creatorName,
    this.assignedToId,
    this.assignedToName,
    required this.createdAt,
    this.updatedAt,
  });
  
  factory Ticket.fromJson(Map<String, dynamic> json) {
    return Ticket(
      id: json['id'],
      title: json['title'],
      description: json['description'],
      status: TicketStatus.values
          .firstWhere((e) => e.name == json['status']),
      priority: TicketPriority.values
          .firstWhere((e) => e.name == json['priority']),
      departmentId: json['departmentId'],
      departmentName: json['departmentName'],
      categoryId: json['categoryId'],
      categoryName: json['categoryName'],
      creatorId: json['creatorId'],
      creatorName: json['creatorName'],
      assignedToId: json['assignedToId'],
      assignedToName: json['assignedToName'],
      createdAt: DateTime.parse(json['createdAt']),
      updatedAt: json['updatedAt'] != null 
          ? DateTime.parse(json['updatedAt']) 
          : null,
    );
  }
}

enum TicketStatus { open, inProgress, pending, resolved, closed }
enum TicketPriority { low, medium, high, critical }
\end{lstlisting}

\newpage

\section{Gerçek Zamanlı İletişim (SignalR)}

\subsection{SignalR Service}

\begin{lstlisting}[language=Dart]
import 'package:signalr_netcore/signalr_client.dart';

class SignalRService {
  HubConnection? _hubConnection;
  
  Future<void> connect(String token) async {
    final serverUrl = 'http://localhost:5000/hubs/ticket';
    
    _hubConnection = HubConnectionBuilder()
        .withUrl(
          serverUrl,
          options: HttpConnectionOptions(
            accessTokenFactory: () => Future.value(token),
          ),
        )
        .withAutomaticReconnect()
        .build();
    
    // Event listeners
    _hubConnection!.on('TicketUpdated', _onTicketUpdated);
    _hubConnection!.on('CommentAdded', _onCommentAdded);
    _hubConnection!.on('TicketAssigned', _onTicketAssigned);
    
    await _hubConnection!.start();
  }
  
  void _onTicketUpdated(List<Object>? args) {
    if (args != null && args.isNotEmpty) {
      final ticketData = args[0] as Map<String, dynamic>;
      print('Ticket updated: ${ticketData['id']}');
      // Notify providers to refresh
    }
  }
  
  void _onCommentAdded(List<Object>? args) {
    if (args != null && args.isNotEmpty) {
      final commentData = args[0] as Map<String, dynamic>;
      print('New comment: ${commentData['content']}');
      // Show notification
    }
  }
  
  Future<void> joinTicket(int ticketId) async {
    await _hubConnection?.invoke('JoinTicket', 
        args: [ticketId]);
  }
  
  Future<void> leaveTicket(int ticketId) async {
    await _hubConnection?.invoke('LeaveTicket', 
        args: [ticketId]);
  }
  
  Future<void> disconnect() async {
    await _hubConnection?.stop();
  }
}
\end{lstlisting}

\newpage

\section{Push Notifications}

\subsection{Notification Service}

\begin{lstlisting}[language=Dart]
import 'package:flutter_local_notifications/
    flutter_local_notifications.dart';

class NotificationService {
  final FlutterLocalNotificationsPlugin _notifications =
      FlutterLocalNotificationsPlugin();
  
  Future<void> initializeLocalNotifications() async {
    const androidSettings = AndroidInitializationSettings(
        '@mipmap/ic_launcher');
    const iosSettings = DarwinInitializationSettings(
      requestAlertPermission: true,
      requestBadgePermission: true,
      requestSoundPermission: true,
    );
    
    const settings = InitializationSettings(
      android: androidSettings,
      iOS: iosSettings,
    );
    
    await _notifications.initialize(
      settings,
      onDidReceiveNotificationResponse: 
          _onNotificationTapped,
    );
  }
  
  Future<void> requestPermissions() async {
    final plugin = _notifications
        .resolvePlatformSpecificImplementation<
            IOSFlutterLocalNotificationsPlugin>();
    
    await plugin?.requestPermissions(
      alert: true,
      badge: true,
      sound: true,
    );
  }
  
  Future<void> showNotification(
    String title,
    String body,
    {int? ticketId}
  ) async {
    const androidDetails = AndroidNotificationDetails(
      'tickly_channel',
      'Tickly Notifications',
      channelDescription: 'Ticket updates and notifications',
      importance: Importance.high,
      priority: Priority.high,
    );
    
    const iosDetails = DarwinNotificationDetails();
    
    const details = NotificationDetails(
      android: androidDetails,
      iOS: iosDetails,
    );
    
    await _notifications.show(
      ticketId ?? 0,
      title,
      body,
      details,
      payload: ticketId?.toString(),
    );
  }
  
  void _onNotificationTapped(
      NotificationResponse response) {
    final payload = response.payload;
    if (payload != null) {
      final ticketId = int.parse(payload);
      // Navigate to ticket detail
    }
  }
}
\end{lstlisting}

\newpage

\section{Kurulum ve Çalıştırma}

\subsection{Geliştirme Ortamı Kurulumu}

\textbf{Gereksinimler:}
\begin{itemize}
    \item Flutter SDK 3.16+
    \item Dart SDK 3.2+
    \item Xcode 14+ (macOS için)
    \item CocoaPods (iOS dependencies için)
    \item iOS Simulator veya fiziksel iPhone
\end{itemize}

\textbf{Flutter SDK Kurulumu:}

\begin{lstlisting}[language=bash]
# macOS için Homebrew ile
brew install flutter

# veya manuel indirme
# https://docs.flutter.dev/get-started/install/macos

# Flutter doctor ile kontrol
flutter doctor
\end{lstlisting}

\subsection{Proje Kurulumu}

\begin{lstlisting}[language=bash]
# 1. Proje klasorune git
cd mobile

# 2. Dependencies yukle
flutter pub get

# 3. iOS dependencies
cd ios
pod install
cd ..

# 4. Cihazlari listele
flutter devices

# 5. iOS Simulator'da calistir
flutter run

# 6. Belirli cihazda calistir
flutter run -d <device-id>
\end{lstlisting}

\subsection{Build ve Release}

\textbf{Development Build:}
\begin{lstlisting}[language=bash]
# Debug mode
flutter run

# Profile mode (performans analizi icin)
flutter run --profile

# Release mode
flutter run --release
\end{lstlisting}

\textbf{iOS App Store Build:}
\begin{lstlisting}[language=bash]
# 1. Build ipa
flutter build ipa

# 2. Xcode ile ac
open ios/Runner.xcworkspace

# 3. Xcode'da Archive > Distribute App
# 4. App Store Connect'e yukle
\end{lstlisting}

\newpage

\section{Tema ve UI Customization}

\subsection{Theme Provider}

\begin{lstlisting}[language=Dart]
class ThemeProvider extends ChangeNotifier {
  ThemeMode _themeMode = ThemeMode.system;
  
  ThemeMode get themeMode => _themeMode;
  
  void toggleTheme() {
    _themeMode = _themeMode == ThemeMode.light 
        ? ThemeMode.dark 
        : ThemeMode.light;
    notifyListeners();
  }
  
  static ThemeData get lightTheme {
    return ThemeData(
      useMaterial3: true,
      colorScheme: ColorScheme.fromSeed(
        seedColor: Colors.blue,
        brightness: Brightness.light,
      ),
      appBarTheme: const AppBarTheme(
        centerTitle: true,
        elevation: 0,
      ),
      cardTheme: CardTheme(
        elevation: 2,
        shape: RoundedRectangleBorder(
          borderRadius: BorderRadius.circular(12),
        ),
      ),
    );
  }
  
  static ThemeData get darkTheme {
    return ThemeData(
      useMaterial3: true,
      colorScheme: ColorScheme.fromSeed(
        seedColor: Colors.blue,
        brightness: Brightness.dark,
      ),
      appBarTheme: const AppBarTheme(
        centerTitle: true,
        elevation: 0,
      ),
      cardTheme: CardTheme(
        elevation: 2,
        shape: RoundedRectangleBorder(
          borderRadius: BorderRadius.circular(12),
        ),
      ),
    );
  }
}
\end{lstlisting}

\newpage

\section{Test Stratejisi}

\subsection{Unit Tests}

\begin{lstlisting}[language=Dart]
import 'package:flutter_test/flutter_test.dart';

void main() {
  group('TicketProvider Tests', () {
    test('Initial state should be empty', () {
      final provider = TicketProvider();
      expect(provider.tickets, isEmpty);
      expect(provider.isLoading, false);
    });
    
    test('Fetch tickets should update state', () async {
      final provider = TicketProvider();
      await provider.fetchTickets();
      
      expect(provider.isLoading, false);
      expect(provider.tickets, isNotEmpty);
    });
  });
  
  group('AuthProvider Tests', () {
    test('Login with valid credentials', () async {
      final provider = AuthProvider();
      final success = await provider.login(
        'test@example.com',
        'password123'
      );
      
      expect(success, true);
      expect(provider.isAuthenticated, true);
      expect(provider.user, isNotNull);
    });
  });
}
\end{lstlisting}

\subsection{Widget Tests}

\begin{lstlisting}[language=Dart]
void main() {
  testWidgets('Login screen has email and password fields',
      (WidgetTester tester) async {
    await tester.pumpWidget(
      const MaterialApp(home: LoginScreen())
    );
    
    expect(find.byType(TextField), findsNWidgets(2));
    expect(find.text('Email'), findsOneWidget);
    expect(find.text('Password'), findsOneWidget);
    expect(find.byType(ElevatedButton), findsOneWidget);
  });
  
  testWidgets('Ticket card displays correct information',
      (WidgetTester tester) async {
    final ticket = Ticket(
      id: 1,
      title: 'Test Ticket',
      description: 'Test Description',
      status: TicketStatus.open,
      priority: TicketPriority.high,
      // ...
    );
    
    await tester.pumpWidget(
      MaterialApp(
        home: Scaffold(
          body: TicketCard(ticket: ticket),
        ),
      ),
    );
    
    expect(find.text('Test Ticket'), findsOneWidget);
    expect(find.text('Test Description'), findsOneWidget);
  });
}
\end{lstlisting}

\newpage

\section{Performans Optimizasyonu}

\subsection{Best Practices}

\begin{enumerate}
    \item \textbf{Lazy Loading:}
    \begin{itemize}
        \item ListView.builder kullanımı
        \item Infinite scroll pagination
        \item Image lazy loading
    \end{itemize}
    
    \item \textbf{State Management:}
    \begin{itemize}
        \item Provider ile efficient state updates
        \item notifyListeners() minimal kullanımı
        \item Selector widget ile specific rebuilds
    \end{itemize}
    
    \item \textbf{API Optimization:}
    \begin{itemize}
        \item Request caching
        \item Debouncing search queries
        \item Parallel requests where possible
    \end{itemize}
    
    \item \textbf{Image Optimization:}
    \begin{itemize}
        \item CachedNetworkImage kullanımı
        \item Image compression
        \item Thumbnail preview
    \end{itemize}
    
    \item \textbf{Build Optimization:}
    \begin{itemize}
        \item const constructors
        \item Widget extraction
        \item RepaintBoundary kullanımı
    \end{itemize}
\end{enumerate}

\newpage

\section{Sonuç ve Gelecek Planları}

\subsection{Mevcut Durum}

Tickly Mobile uygulaması, iOS platformu için Flutter kullanılarak başarıyla geliştirilmiştir. Uygulama aşağıdaki temel özellikleri sunar:

\begin{itemize}
    \item ✓ JWT tabanlı güvenli kimlik doğrulama
    \item ✓ Ticket oluşturma, listeleme, detay görüntüleme
    \item ✓ Yorum sistemi
    \item ✓ Dosya ve fotoğraf yükleme
    \item ✓ Real-time SignalR entegrasyonu
    \item ✓ Push notification desteği
    \item ✓ Dark mode
    \item ✓ Dashboard ve istatistikler
\end{itemize}

\subsection{Gelecek Geliştirmeler}

\textbf{Kısa Vadeli (1-2 ay):}
\begin{itemize}
    \item Android platform desteği
    \item Offline-first architecture (local database)
    \item Biometric authentication (Face ID / Touch ID)
    \item Rich text editor (yorumlar için)
    \item File preview in-app
\end{itemize}

\textbf{Orta Vadeli (2-4 ay):}
\begin{itemize}
    \item Voice-to-text ticket oluşturma
    \item QR code ile ticket açma
    \item Advanced filtering ve sorting
    \item Ticket templates
    \item Export functionality (PDF)
\end{itemize}

\textbf{Uzun Vadeli (4-6 ay):}
\begin{itemize}
    \item iPad optimization (tablet layout)
    \item Apple Watch companion app
    \item Siri Shortcuts integration
    \item AR features (problem visualization)
    \item ML-based smart suggestions
\end{itemize}

\subsection{Teknik İyileştirmeler}

\begin{itemize}
    \item Comprehensive test coverage (\%80+)
    \item CI/CD pipeline (GitHub Actions + Fastlane)
    \item Performance monitoring (Firebase Performance)
    \item Crash reporting (Firebase Crashlytics)
    \item Analytics integration (Firebase Analytics)
    \item Code generation (freezed, json\_serializable)
\end{itemize}

\vspace{1cm}

\begin{center}
\textit{--- Doküman Sonu ---}
\end{center}

\end{document}
